\section{失败解释、科学边界与AI使用}
\label{cr:discussion}
\subsection{工程恢复不会增加科学样本}
R50主线经历四次原始工程失败，其有效确定性重试分别由r2、r3和r4版本接续。
原始失败的退出码未被观测时仍保留为空值；后续收集成功或重试成功不能改写原始退出码。
四个alias只把真实通过的归档放回原逻辑槽，主分母仍是64个独立事件及128个配对响应。
正式统计不读取原FAILED的电学输出来选择更有利的结果，全部历史成本留在唯一账本。
原100~$\Omega$预充失败与独立probe另属历史阶段，也不加入R50主矩阵。

工程完成与科学有效性需分别判断。完整归档和退出0只能证明执行链满足工程条件；
初态、预充、积分、守恒、真实终点与匹配参考仍须通过科学审计。
旧100~$\Omega$预充误差超限时，即使最终裕量看似安全也不能纳入主线。
50~$\Omega$修订改变了外部电路，依靠新版本、新参考和新响应建立证据，原阈值没有放宽。

\subsection{连续量、读出判据与数值不确定性}
同方向的SN扰动对存储0与存储1可能有不同影响。
式\eqref{cr:metrics}用位值符号定义裕量，因而$\Delta M$可以为负，表示相对无粒子对照的裕量增加。
生成电荷、端口收集电荷和存储电容净变化不能互相替代；复合、不同端口收集及电容耦合共同作用。
因此不能以$Q_{\mathrm{dep}}/C$代替实际瞬态求解。

本次边界验证按完整128响应的预定距离排序选择两个目标，不以输出是否有利改变目标。
连续变化通过容差，不等于全部离散类别已确定。未选中但仍距某阈值不超过0.5~mV的响应，
若没有规定的两项匹配细化，仍按近界不确定处理。
原始分类与最终有效分类同时保留；临界和类别跨界不会自动抹去有限的连续目标，
连续数值失败则传播到有效评分标记。

每个位值/尺寸/粒子/方向分层只有四个独立事件。零次明确误读不能证明概率为零，
未知和临界状态需要用固定分母及最佳/最差包络解释。配对两位、工程重试和扩展不能把N翻倍。
本研究没有真实感放锁存、恢复与长期保持测试，固定4~ns采样的判据误读不能解释为持久位翻转。

\subsection{预测验证的含义}
模型先于测试TCAD提交冻结，测试仅用于评价固定预测。域内留出和尺寸外推分别给出
MAE、RMSE、最大绝对误差及基准比较，不将训练LOCO误差当作独立测试性能。
模型的输入包含真实Geant4沉积量，其用途是给定源项及基线条件后的TCAD响应近似，
并未替代全部输运或基准仿真。
本次完整测试评价扩大了已验证范围，但范围仍限于指定四栅长、两粒子、单一5~MeV能量和两个方向。
不能由固定集合表现推断任意粒子谱、入射位置或新工艺节点。

\subsection{受控初态、二维模型和材料栈限制}
存储0/1由理想钳位及准静态初态建立，随后释放钳位并检查粒子前电压。
这一设计控制初始条件，不能说明实际访问管能够在同样时间内写入，亦未校准写入能耗或保持可靠性。
二维平面器件、1~$\mu$m有效厚度、理想集总电容和固定掺杂均是探索设定。
数值一致性只能检验当前方程和离散实现，不能发现共同的模型系统偏差。

Geant4采用硅/真空几何与电磁option4。质子结果限于所用电磁配置的电离响应，
不能写为已经验证核反应次级粒子或核反冲主导的SEE。
4/6~nm氧化层结果使用同一硅入口事件，只是TCAD参数敏感性；未重算变化材料栈中的完整输运。
总量等效直径迹也没有保留完整三维逐步能量沉积的空间和时间分布。

\subsection{AI作用与课程责任}
AI协助文献组织、脚本开发、失败取证、证据核验、图表和文稿编写；固定代码承担参数映射、
求解顺序、阈值判断和评分。没有开展人工操作对照实验，不能报告AI加速倍数或因果效率优势。
工程流水线完成不说明学生已独立理解或操作；小组人数、真实贡献、教师审核、资源交接、
组内试讲、现场演示、答辩与成绩均保留待人工记录。
后续器件研究需要有出处的工艺结构和电学数据校准，课程验收则需要真实参与者完成相应教学活动。
